%  LaTeX support: latex@mdpi.com 
%  For support, please attach all files needed for compiling as well as the log file, and specify your operating system, LaTeX version, and LaTeX editor.

%=================================================================
\documentclass[journal,article,submit,pdftex,moreauthors]{Definitions/mdpi} 
%\documentclass[preprints,article,submit,pdftex,moreauthors]{Definitions/mdpi} 
% For posting an early version of this manuscript as a preprint, you may use "preprints" as the journal. Changing "submit" to "accept" before posting will remove line numbers.

% Below journals will use APA reference format:
% admsci, aieduc, behavsci, businesses, econometrics, economies, education, ejihpe, famsci, games, humans, ijcs, ijfs, jintelligence, journalmedia, jrfm, jsam, languages, peacestud, psycholint, publications, tourismhosp, youth

% Below journals will use Chicago reference format:
% arts, genealogy, histories, humanities, laws, literature, religions, risks, socsci

%--------------------
% Class Options:
%--------------------
%----------
% journal
%----------
% Choose between the following MDPI journals:
% accountaudit, acoustics, actuators, addictions, adhesives, admsci, adolescents, aerobiology, aerospace, agriculture, agriengineering, agrochemicals, agronomy, ai, aichem, aieduc, aieng, aimater, aimed, aipa, air, aisens, algorithms, allergies, alloys, amh, analog, analytica, analytics, anatomia, anesthres, animals, antibiotics, antibodies, antioxidants, applbiosci, appliedchem, appliedmath, appliedphys, applmech, applmicrobiol, applnano, applsci, aquacj, architecture, arm, arthropoda, arts, asc, asi, astronautics, astronomy, atmosphere, atoms, audiolres, automation, axioms, bacteria, batteries, bdcc, behavsci, beverages, biochem, bioengineering, biologics, biology, biomass, biomechanics, biomed, biomedicines, biomedinformatics, biomimetics, biomolecules, biophysica, bioresourbioprod, biosensors, biosphere, biotech, birds, blockchains, bloods, blsf, brainsci, breath, buildings, businesses, cancers, carbon, cardio, cardiogenetics, cardiovascmed, catalysts, cells, ceramics, challenges, chemengineering, chemistry, chemosensors, chemproc, children, chips, cimb, civileng, cleantechnol, climate, clinbioenerg, clinpract, clockssleep, cmd, cmtr, coasts, coatings, colloids, colorants, commodities, complexities, complications, compounds, computation, computers, condensedmatter, conservation, constrmater, cosmetics, covid, crops, cryo, cryptography, crystals, csmf, ctn, culture, curroncol, cyber, dairy, data, ddc, dentistry, dermato, dermatopathology, designs, devices, dhi, diabetology, diagnostics, dietetics, digital, disabilities, diseases, diversity, dna, drones, dynamics, earth, ebj, ecm, ecologies, econometrics, economies, edm, education, eesp, ejihpe, electricity, electrochem, electronicmat, electronics, encyclopedia, endocrines, energies, eng, engproc, entomology, entropy, environments, environremediat, epidemiologia, epigenomes, esa, est, famsci, fermentation, fibers, fintech, fire, fishes, fluids, foods, forecasting, forensicsci, forests, fossstud, foundations, fractalfract, fuels, future, futureinternet, futurepharmacol, futurephys, futuretransp, galaxies, games, gases, gastroent, gastrointestdisord, gastronomy, gels, genealogy, genes, geographies, geohazards, geomatics, geometry, geosciences, geotechnics, geriatrics, germs, glacies, grasses, green, greenhealth, gucdd, hardware, hazardousmatters, healthcare, hearts, hemato, hematolrep, hep, heritage, higheredu, highthroughput, histories, horticulturae, hospitals, humanities, humans, hydrobiology, hydrogen, hydrology, hydropower, hygiene, idr, iic, ijcs, ijem, ijerph, ijfs, ijgi, ijmd, ijms, ijns, ijom, ijpb, ijt, ijtm, ijtpp, ime, immuno, informatics, information, infrastructures, inorganics, insects, instruments, inventions, iot, j, jaestheticmed, jal, jcdd, jcm, jcp, jcrm, jcs, jcto, jdad, jdb, jdream, jemr, jeta, jfb, jfmk, jgbg, jgg, jimaging, jintelligence, jlpea, jmahp, jmmp, jmms, jmp, jmse, jne, jnt, jof, joi, joitmc, joma, jop, joptm, jor, journalmedia, jox, jpbi, jphytomed, jpm, jrfm, jsam, jsan, jtaer, jvd, jzbg, kidneydial, kinasesphosphatases, knowledge, labmed, laboratories, lae, land, languages, laws, life, lights, limnolrev, lipidology, liquids, literature, livers, logics, logistics, lubricants, lymphatics, machines, macromol, magnetism, magnetochemistry, make, marinedrugs, materials, materproc, mathematics, mca, measurements, medicina, medicines, medsci, membranes, merits, metabolites, metals, meteorology, methane, metrics, metrology, micro, microarrays, microbiolres, microelectronics, micromachines, microorganisms, microplastics, microwave, minerals, mining, mmphys, modelling, molbank, molecules, mps, msf, mti, multimedia, muscles, nanoenergyadv, nanomanufacturing, nanomaterials, ncrna, ndt, network, neuroglia, neuroimaging, neurolint, neurosci, nitrogen, notspecified, nursrep, nutraceuticals, nutrients, obesities, occuphealth, oceans, ohbm, onco, optics, oral, organics, organoids, osteology, oxygen, pandemics, parasites, parasitologia, particles, pathogens, pathophysiology, peacestud, pediatrrep, pets, pharmaceuticals, pharmaceutics, pharmacoepidemiology, pharmacy, philosophies, photochem, photonics, phycology, physchem, physics, physiologia, plants, plasma, platforms, pollutants, polymers, polysaccharides, populations, poultry, powders, precipitation, precisoncol, preprints, proceedings, processes, prosthesis, proteomes, psf, psychiatryint, psychoactives, psycholint, publications, purification, quantumrep, quaternary, qubs, radiation, rdt, reactions, realestate, receptors, recycling, regeneration, religions, remotesensing, reports, reprodmed, resources, rheumato, risks, rjpm, robotics, rsee, ruminants, safety, sci, scipharm, sclerosis, seeds, sensors, separations, sexes, shi, signals, sinusitis, siuj, skins, smartcities, sna, societies, socsci, software, soilsystems, solar, solids, spectroscj, sports, standards, stats, std, stratsediment, stresses, surfaces, surgeries, suschem, sustainability, symmetry, synbio, systems, tae, targets, taxonomy, technologies, telecom, test, textiles, thalassrep, therapeutics, thermo, timespace, tomography, tourismhosp, toxics, toxins, tph, transplantology, transportation, traumacare, traumas, tri, tropicalmed, universe, urbansci, uro, vaccines, vehicles, venereology, vetsci, vibration, virtualworlds, viruses, vision, waste, water, welding, wem, wevj, wild, wind, women, world, youth, zoonoticdis

%---------
% article
%---------
% The default type of manuscript is "article", but can be replaced by: 
% abstract, addendum, article, benchmark, book, bookreview, briefcommunication, briefreport, casereport, changes, clinicopathologicalchallenge, comment, commentary, communication, conceptpaper, conferenceproceedings, correction, conferencereport, creative, datadescriptor, discussion, entry, expressionofconcern, extendedabstract, editorial, essay, erratum, fieldguide, hypothesis, interestingimages, letter, meetingreport, monograph, newbookreceived, obituary, opinion, proceedingpaper, projectreport, reply, retraction, review, perspective, protocol, shortnote, studyprotocol, supfile, systematicreview, technicalnote, viewpoint, guidelines, registeredreport, tutorial,  giantsinurology, urologyaroundtheworld
% supfile = supplementary materials

%----------
% submit
%----------
% The class option "submit" will be changed to "accept" by the Editorial Office when the paper is accepted. This will only make changes to the frontpage (e.g., the logo of the journal will get visible), the headings, and the copyright information. Also, line numbering will be removed. Journal info and pagination for accepted papers will also be assigned by the Editorial Office.

%------------------
% moreauthors
%------------------
% If there is only one author the class option oneauthor should be used. Otherwise use the class option moreauthors.

%---------
% pdftex
%---------
% The option pdftex is for use with pdfLaTeX. Remove "pdftex" for (1) compiling with LaTeX & dvi2pdf (if eps figures are used) or for (2) compiling with XeLaTeX.

%=================================================================
% MDPI internal commands - do not modify
\firstpage{1} 
\makeatletter 
\setcounter{page}{\@firstpage} 
\makeatother
\pubvolume{1}
\issuenum{1}
\articlenumber{0}
\pubyear{2026}
\copyrightyear{2026}
%\externaleditor{Firstname Lastname} % More than 1 editor, please add `` and '' before the last editor name
\datereceived{ } 
\daterevised{ } % Comment out if no revised date
\dateaccepted{ } 
\datepublished{ } 
%\datecorrected{} % For corrected papers: "Corrected: XXX" date in the original paper.
%\dateretracted{} % For retracted papers: "Retracted: XXX" date in the original paper.
%\doinum{} % Used for some special journals, like molbank
%\pdfoutput=1 % Uncommented for upload to arXiv.org
%\CorrStatement{yes}  % For updates
%\longauthorlist{yes} % For many authors that exceed the left citation part
%\IsAssociation{yes} % For association journals

%=================================================================
% Add packages and commands here. The following packages are loaded in our class file: fontenc, inputenc, calc, indentfirst, fancyhdr, graphicx, epstopdf, lastpage, ifthen, float, amsmath, amssymb, lineno, setspace, enumitem, mathpazo, booktabs, titlesec, etoolbox, tabto, xcolor, colortbl, soul, multirow, microtype, tikz, totcount, changepage, attrib, upgreek, array, tabularx, pbox, ragged2e, tocloft, marginnote, marginfix, enotez, amsthm, natbib, hyperref, cleveref, scrextend, url, geometry, newfloat, caption, draftwatermark, seqsplit
% cleveref: load \crefname definitions after \begin{document}

%=================================================================
% Please use the following mathematics environments: Theorem, Lemma, Corollary, Proposition, Characterization, Property, Problem, Example, ExamplesandDefinitions, Hypothesis, Remark, Definition, Notation, Assumption
%% For proofs, please use the proof environment (the amsthm package is loaded by the MDPI class).

%=================================================================
% Full title of the paper (Capitalized)
\Title{Sectoral Prioritization in Turkey’s Socio-Economic Production System: Integrating Production-Network Centrality and CGE Impact Analysis}

% Author Orchid ID: enter ID or remove command
\newcommand{\orcidauthorA}{0000-0000-0000-000X} % Add \orcidA{} behind the author's name
%\newcommand{\orcidauthorB}{0000-0000-0000-000X} % Add \orcidB{} behind the author's name

% Authors, for the paper (add full first names)
\Author{Firstname Lastname $^{1}$\orcidA{}, Firstname Lastname $^{2}$ and Firstname Lastname $^{2,}$*}

%\longauthorlist{yes}

% MDPI internal command: Authors, for metadata in PDF
\AuthorNames{Firstname Lastname, Firstname Lastname and Firstname Lastname}

%\longauthorlist{yes}

% Affiliations / Addresses (Add [1] after \address if there is only one affiliation.)
\address{%
$^{1}$ \quad Affiliation 1; e-mail@e-mail.com\\
$^{2}$ \quad Affiliation 2; e-mail@e-mail.com}

% Contact information of the corresponding author
\corres{Correspondence: e-mail@e-mail.com; Tel.: (optional; include country code; if there are multiple corresponding authors, add author initials) +xx-xxxx-xxx-xxxx (F.L.)}

% Current address and/or shared authorship
%\firstnote{Current address: Affiliation.}  
% Current address should not be the same as any items in the Affiliation section.

%\secondnote{These authors contributed equally to this work.}
% The commands \thirdnote{} till \eighthnote{} are available for further notes.

%\simplesumm{} % Simple summary

%\conference{} % An extended version of a conference paper

% Abstract (Do not insert blank lines, i.e. \\) 
\abstract{Sectoral prioritization can conflate a sector’s structural position in the production network with its economy-wide policy impact. This study tests that assumption using Turkey’s 2023 production network and a computable general equilibrium model calibrated to a Social Accounting Matrix. For 62 sectors, it compares the Structural Network Importance Index (SNII), which captures network position, with the Economy-wide Impact Performance Index (EIPI), which summarizes productivity-shock effects on growth, welfare, external balance, and prices after controlling for sector size. Results show that centrality alone does not imply a large or balanced general-equilibrium impact. Network indicators contain useful but limited information about production and spillover effects, while price, foreign-trade, and distributional outcomes depend more strongly on sector size, imported-input dependence, and price adjustment. We therefore interpret centrality as complementary structural information rather than a substitute for general-equilibrium analysis. The framework combines a two-dimensional sector map with Pareto prioritization, balance constraints, and specification-robustness checks to identify a robust structural core and policy-sensitive candidates. The findings support multidimensional, objective-sensitive industrial policy rather than a one-dimensional list of winning sectors.}

% Keywords
\keyword{socio-economic systems; systems analysis; production networks; input-output analysis; computable general equilibrium; network centrality; sectoral prioritization; industrial policy; decision support; Turkey} 

% The fields PACS, MSC, and JEL may be left empty or commented out if not applicable
%\PACS{J0101}
%\MSC{}
%\JEL{}

%%%%%%%%%%%%%%%%%%%%%%%%%%%%%%%%%%%%%%%%%%
% Only for the journal Diversity
%\LSID{\url{http://}}

%%%%%%%%%%%%%%%%%%%%%%%%%%%%%%%%%%%%%%%%%%
% Only for the journal Applied Sciences
%\featuredapplication{Authors are encouraged to provide a concise description of the specific application or a potential application of the work. This section is not mandatory.}
%%%%%%%%%%%%%%%%%%%%%%%%%%%%%%%%%%%%%%%%%%

%%%%%%%%%%%%%%%%%%%%%%%%%%%%%%%%%%%%%%%%%%
% Only for the journal Data
%\dataset{DOI number or link to the deposited data set if the data set is published separately. If the data set shall be published as a supplement to this paper, this field will be filled by the journal editors. In this case, please submit the data set as a supplement.}
%\datasetlicense{License under which the data set is made available (CC0, CC-BY, CC-BY-SA, CC-BY-NC, etc.)}

%%%%%%%%%%%%%%%%%%%%%%%%%%%%%%%%%%%%%%%%%%
% Only for the journal BioTech, Fishes, Neuroimaging and Toxins
%\keycontribution{The breakthroughs or highlights of the manuscript. Authors can write one or two sentences to describe the most important part of the paper.}

%%%%%%%%%%%%%%%%%%%%%%%%%%%%%%%%%%%%%%%%%%
% Only for the journal Encyclopedia
%\encyclopediadef{For entry manuscripts only: please provide a brief overview of the entry title instead of an abstract.}

%%%%%%%%%%%%%%%%%%%%%%%%%%%%%%%%%%%%%%%%%%
% Different journals have different requirements. Please check the specific journal guidelines in the "Instructions for Authors" on the journal's official website.
%\addhighlights{yes}
%\renewcommand{\addhighlights}{%
%
%\noindent The goal is to increase the discoverability and readability of the article via search engines and other scholars. Highlights should not be a copy of the abstract, but a simple text allowing the reader to quickly and simplified find out what the article is about and what can be cited from it. Each of these parts should be devoted up to 2~bullet points.\vspace{3pt}\\
%\textbf{What are the main findings?}
% \begin{itemize}[labelsep=2.5mm,topsep=-3pt]
% \item First bullet.
% \item Second bullet.
% \end{itemize}\vspace{3pt}
%\textbf{What are the implications of the main findings?}
% \begin{itemize}[labelsep=2.5mm,topsep=-3pt]
% \item First bullet.
% \item Second bullet.
% \end{itemize}
%}

%%%%%%%%%%%%%%%%%%%%%%%%%%%%%%%%%%%%%%%%%%
\begin{document}

%%%%%%%%%%%%%%%%%%%%%%%%%%%%%%%%%%%%%%%%%%

\section{Introduction}

The introduction should briefly place the study in a broad context and highlight why it is important. It should define the purpose of the work and its significance. The current state of the research field should be reviewed carefully and key publications cited. Please highlight controversial and diverging hypotheses when necessary. Finally, briefly mention the main aim of the work and highlight the principal conclusions. As far as possible, please keep the introduction comprehensible to scientists outside your particular field of research. Citing a journal paper \citep{}.  Now citing a book reference \citep{ref-book1,ref-book2} or other reference types \citep{ref-unpublish,ref-url}. Please use the command \citep{ref-proceeding,ref-thesis} for the following MDPI journals, which use author--date citation: Administrative Sciences, Arts, Behavioral Sciences, Businesses, Econometrics, Economies, Education Sciences, European Journal of Investigation in Health, Psychology and Education, Games, Genealogy, Histories, Humanities, Humans, IJFS, Journal of Intelligence, Journalism and Media, JRFM, Languages, Laws, Literature, Psychology International, Publications, Religions, Risks, Social Sciences, Tourism and Hospitality, Youth. 

%%%%%%%%%%%%%%%%%%%%%%%%%%%%%%%%%%%%%%%%%%
\section{Materials and Methods}

Materials and Methods should be described with sufficient details to allow others to replicate and build on published results. Please note that publication of your manuscript implicates that you must make all materials, data, computer code, and protocols associated with the publication available to readers. Please disclose at the submission stage any restrictions on the availability of materials or information. New methods and protocols should be described in detail while well-established methods can be briefly described and appropriately cited.

Research manuscripts reporting large datasets that are deposited in a publicly avail-able database should specify where the data have been deposited and provide the relevant accession numbers. If the accession numbers have not yet been obtained at the time of submission, please state that they will be provided during review. They must be provided prior to publication.

Interventionary studies involving animals or humans, and other studies require ethical approval must list the authority that provided approval and the corresponding ethical approval code.

In this section, where applicable, authors are required to disclose details of how gen-erative artificial intelligence (GenAI) has been used in this paper (e.g., to generate text, data, or graphics, or to assist in study design, data collection, analysis, or interpretation). The use of GenAI for superficial text editing (e.g., grammar, spelling, punctuation, and formatting) does not need to be declared.

\begin{quote}
This is an example of a quote.
\end{quote}

%%%%%%%%%%%%%%%%%%%%%%%%%%%%%%%%%%%%%%%%%%
\section{Results}

This section may be divided by subheadings. It should provide a concise and precise description of the experimental results, their interpretation as well as the experimental conclusions that can be drawn.
\subsection{Subsection}
\subsubsection{Subsubsection}

Bulleted lists look like this:
\begin{itemize}
\item	First bullet;
\item	Second bullet;
\item	Third bullet.
\end{itemize}

Numbered lists can be added as follows:
\begin{enumerate}
\item	First item; 
\item	Second item;
\item	Third item.
\end{enumerate}

The text continues here.

\subsection{Figures, Tables and Schemes}

All figures and tables should be cited in the main text as Figure~\ref{fig1}, Table~\ref{tab1}, etc.

\begin{figure}[H]
%\isPreprints{\centering}{} % Only used for preprints
\includegraphics[width=4.0 cm]{Definitions/logo-mdpi}
\caption{This is a figure. Schemes follow the same formatting.\label{fig1}}
\end{figure}   
\unskip

\begin{table}[H] 
%\small % Change table font size
\caption{This is a table caption. Tables should be placed in the main text near to the first time they are~cited.\label{tab1}}
%\isPreprints{\centering}{} % Only used for preprints
\begin{tabularx}{\textwidth}{CCC}
\toprule
\textbf{Title 1}	& \textbf{Title 2}	& \textbf{Title 3}\\
\midrule
Entry 1		& Data			& Data\\
Entry 2		& Data			& Data \textsuperscript{1}\\
\bottomrule
\end{tabularx}

\noindent{\footnotesize{\textsuperscript{1} Tables may have a footer.}}
\end{table}

The text continues here (Figure~\ref{fig2} and Table~\ref{tab2}).

% Example of a figure that spans the whole page width and with subfigures. The same concept works for tables, too.
\begin{figure}[H]
%\isPreprints{}{% This command is only used for ``preprints''.
\begin{adjustwidth}{-\extralength}{0cm}
\centering
%} % If the paper is ``preprints'', please uncomment this parenthesis.
\subfloat[\centering]{\includegraphics[width=7.0cm]{Definitions/logo-mdpi}}
%\hfill
\subfloat[\centering]{\includegraphics[width=7.0cm]{Definitions/logo-mdpi}}\\
\subfloat[\centering]{\includegraphics[width=7.0cm]{Definitions/logo-mdpi}}
%\hfill
\subfloat[\centering]{\includegraphics[width=7.0cm]{Definitions/logo-mdpi}}
%\isPreprints{}{% This command is only used for ``preprints''.
\end{adjustwidth}
%} % If the paper is ``preprints'', please uncomment this parenthesis.
\caption{This is a wide figure. Schemes follow the same formatting. If there are multiple panels, they should be listed as: (\textbf{a}) Description of what is contained in the first panel. (\textbf{b}) Description of what is contained in the second panel. (\textbf{c}) Description of what is contained in the third panel. (\textbf{d}) Description of what is contained in the fourth panel. Figures should be placed in the main text near to the first time they are cited. A caption on a single line should be centered.\label{fig2}}
\end{figure} 

\begin{table}[H]
\caption{This is a wide table.\label{tab2}}
%\isPreprints{\centering}{% This command is only used for ``preprints''.
	\begin{adjustwidth}{-\extralength}{0cm}
%} % If the paper is ``preprints'', please uncomment this parenthesis.
%\isPreprints{\begin{tabularx}{\textwidth}{CCCC}}{% This command is only used for ``preprints''.
		\begin{tabularx}{\fulllength}{CCCC}
%} % If the paper is ``preprints'', please uncomment this parenthesis.
			\toprule
			\textbf{Title 1}	& \textbf{Title 2}	& \textbf{Title 3}     & \textbf{Title 4}\\
			\midrule
\multirow[m]{3}{*}{Entry 1 *}	& Data			& Data			& Data\\
			  	                   & Data			& Data			& Data\\
			             	      & Data			& Data			& Data\\
                   \midrule
\multirow[m]{3}{*}{Entry 2}    & Data			& Data			& Data\\
			  	                  & Data			& Data			& Data\\
			             	     & Data			& Data			& Data\\
			\bottomrule
		\end{tabularx}
%		\isPreprints{}{% This command is only used for ``preprints''.
	\end{adjustwidth}
%} % If the paper is ``preprints'', please uncomment this parenthesis.
	\noindent{\footnotesize{* Tables may have a footer.}}
\end{table}

%\begin{listing}[H]
%\caption{Title of the listing}
%\rule{\columnwidth}{1pt}
%\raggedright Text of the listing. In font size footnotesize, small, or normalsize. Preferred format: left aligned and single spaced. Preferred border format: top border line and bottom border line.
%\rule{\columnwidth}{1pt}
%\end{listing}

Text.

Text.

\subsection{Formatting of Mathematical Components}

This is the example 1 of equation:
\begin{linenomath}
\begin{equation}
a = 1,
\end{equation}
\end{linenomath}
the text following an equation need not be a new paragraph. Please punctuate equations as regular text.
%% If the documentclass option "submit" is chosen, please insert a blank line before and after any math environment (equation and eqnarray environments). This ensures correct linenumbering. The blank line should be removed when the documentclass option is changed to "accept" because the text following an equation should not be a new paragraph.

This is the example 2 of equation:
%\isPreprints{}{% This command is only used for ``preprints''.
\begin{adjustwidth}{-\extralength}{0cm}
%} % If the paper is ``preprints'', please uncomment this parenthesis.
\begin{equation}
a = b + c + d + e + f + g + h + i + j + k + l + m + n + o + p + q + r + s + t + u + v + w + x + y + z
\end{equation}
%\isPreprints{}{% This command is only used for ``preprints''.
\end{adjustwidth}
%} % If the paper is ``preprints'', please uncomment this parenthesis.

%% Example of a page in landscape format (with table and table footnote).
%\startlandscape
%\begin{table}[H] %% Table in wide page
%%\isPreprints{\centering}{} % This command is only used for ``preprints''.
%\caption{This is a very wide table.\label{tab3}}
%	\begin{tabularx}{\textwidth}{CCCC}
%		\toprule
%		\textbf{Title 1}	& \textbf{Title 2}	& \textbf{Title 3}	& \textbf{Title 4}\\
%		\midrule
%		Entry 1		& Data			& Data			& This cell has some longer content that runs over two lines.\\
%		Entry 2		& Data			& Data			& Data\textsuperscript{1}\\
%		\bottomrule
%	\end{tabularx}
%%\isPreprints{}{% This command is only used for ``preprints''.
%	\begin{adjustwidth}{+\extralength}{0cm}
%%} % If the paper is ``preprints'', please uncomment this parenthesis.
%		\noindent\footnotesize{\textsuperscript{1} This is a table footnote.}
%%\isPreprints{}{% This command is only used for ``preprints''.
%	\end{adjustwidth}
%%} % If the paper is ``preprints'', please uncomment this parenthesis.
%\end{table}
%\finishlandscape

Please punctuate equations as regular text. Theorem-type environments (including propositions, lemmas, corollaries etc.) can be formatted as follows:
%% Example of a theorem:
\begin{Theorem}
Example text of a theorem.
\end{Theorem}

The text continues here. Proofs must be formatted as follows:

%% Example of a proof:
\begin{proof}[Proof of Theorem 1]
Text of the proof. Note that the phrase ``of Theorem 1'' is optional if it is clear which theorem is being referred to.
\end{proof}
The text continues here.

%%%%%%%%%%%%%%%%%%%%%%%%%%%%%%%%%%%%%%%%%%
\section{Discussion}

Authors should discuss the results and how they can be interpreted from the perspective of previous studies and of the working hypotheses. The findings and their implications should be discussed in the broadest context possible. Future research directions may also be highlighted.

%%%%%%%%%%%%%%%%%%%%%%%%%%%%%%%%%%%%%%%%%%
\section{Conclusions}

This section is not mandatory, but can be added to the manuscript if the discussion is unusually long or complex.

%%%%%%%%%%%%%%%%%%%%%%%%%%%%%%%%%%%%%%%%%%
\section{Patents}

This section is not mandatory, but may be added if there are patents resulting from the work reported in this manuscript.

%%%%%%%%%%%%%%%%%%%%%%%%%%%%%%%%%%%%%%%%%%
\vspace{6pt} 

%%%%%%%%%%%%%%%%%%%%%%%%%%%%%%%%%%%%%%%%%%
%% optional
%\supplementary{The following supporting information can be downloaded at:  \linksupplementary{s1}, Figure S1: title; Table S1: title; Video S1: title.}

% Only for journal Methods and Protocols:
% If you wish to submit a video article, please do so with any other supplementary material.
% \supplementary{The following supporting information can be downloaded at: \linksupplementary{s1}, Figure S1: title; Table S1: title; Video S1: title. A supporting video article is available at doi: link.}

% Only used for preprtints:
% \supplementary{The following supporting information can be downloaded at the website of this paper posted on \href{https://www.preprints.org/}{Preprints.org}.}

% Only for journal Hardware:
% If you wish to submit a video article, please do so with any other supplementary material.
% \supplementary{The following supporting information can be downloaded at: \linksupplementary{s1}, Figure S1: title; Table S1: title; Video S1: title.\vspace{6pt}\\
%\begin{tabularx}{\textwidth}{lll}
%\toprule
%\textbf{Name} & \textbf{Type} & \textbf{Description} \\
%\midrule
%S1 & Python script (.py) & Script of python source code used in XX \\
%S2 & Text (.txt) & Script of modelling code used to make Figure X \\
%S3 & Text (.txt) & Raw data from experiment X \\
%S4 & Video (.mp4) & Video demonstrating the hardware in use \\
%... & ... & ... \\
%\bottomrule
%\end{tabularx}
%}

%%%%%%%%%%%%%%%%%%%%%%%%%%%%%%%%%%%%%%%%%%
\authorcontributions{For research articles with several authors, a short paragraph specifying their individual contributions must be provided. The following statements should be used ``Conceptualization, X.X. and Y.Y.; methodology, X.X.; software, X.X.; validation, X.X., Y.Y. and Z.Z.; formal analysis, X.X.; investigation, X.X.; resources, X.X.; data curation, X.X.; writing---original draft preparation, X.X.; writing---review and editing, X.X.; visualization, X.X.; supervision, X.X.; project administration, X.X.; funding acquisition, Y.Y. All authors have read and agreed to the published version of the manuscript.'', please turn to the  \href{http://img.mdpi.org/data/contributor-role-instruction.pdf}{CRediT taxonomy} for the term explanation. Authorship must be limited to those who have contributed substantially to the work~reported.}

\funding{Please add: ``This research received no external funding'' or ``This research was funded by NAME OF FUNDER grant number XXX.'' and  and ``The APC was funded by XXX''. Check carefully that the details given are accurate and use the standard spelling of funding agency names at \url{https://search.crossref.org/funding}, any errors may affect your future funding.}

\institutionalreview{In this section, you should add the Institutional Review Board Statement and approval number, if relevant to your study. You might choose to exclude this statement if the study did not require ethical approval. Please note that the Editorial Office might ask you for further information. Please add “The study was conducted in accordance with the Declaration of Helsinki, and approved by the Institutional Review Board (or Ethics Committee) of NAME OF INSTITUTE (protocol code XXX and date of approval).” for studies involving humans. OR “The animal study protocol was approved by the Institutional Review Board (or Ethics Committee) of NAME OF INSTITUTE (protocol code XXX and date of approval).” for studies involving animals. OR “Ethical review and approval were waived for this study due to REASON (please provide a detailed justification).” OR “Not applicable” for studies not involving humans or animals.}

\informedconsent{Any research article describing a study involving humans should contain this statement. Please add ``Informed consent was obtained from all subjects involved in the study.'' OR ``Patient consent was waived due to REASON (please provide a detailed justification).'' OR ``Not applicable'' for studies not involving humans. You might also choose to exclude this statement if the study did not involve humans.

Written informed consent for publication must be obtained from participating patients who can be identified (including by the patients themselves). Please state ``Written informed consent has been obtained from the patient(s) to publish this paper'' if applicable.}

\dataavailability{We encourage all authors of articles published in MDPI journals to share their research data. In this section, please provide details regarding where data supporting reported results can be found, including links to publicly archived datasets analyzed or generated during the study. Where no new data were created, or where data is unavailable due to privacy or ethical restrictions, a statement is still required. Suggested Data Availability Statements are available in section ``MDPI Research Data Policies'' at \url{https://www.mdpi.com/ethics}.} 

%\durcstatement{Current research is limited to the [please insert a specific academic field, e.g., XXX], which is beneficial [share benefits and/or primary use] and does not pose a threat to public health or national security. Authors acknowledge the dual-use potential of the research involving xxx and confirm that all necessary precautions have been taken to prevent potential misuse. As an ethical responsibility, authors strictly adhere to relevant national and international laws about DURC. Authors advocate for responsible deployment, ethical considerations, regulatory compliance, and transparent reporting to mitigate misuse risks and foster beneficial outcomes.}

% Only for journal Nursing Reports
%\publicinvolvement{Please describe how the public (patients, consumers, carers) were involved in the research. Consider reporting against the GRIPP2 (Guidance for Reporting Involvement of Patients and the Public) checklist. If the public were not involved in any aspect of the research add: ``No public involvement in any aspect of this research''.}
%
%% Only for journal Nursing Reports
%\guidelinesstandards{Please add a statement indicating which reporting guideline was used when drafting the report. For example, ``This manuscript was drafted against the XXX (the full name of reporting guidelines and citation) for XXX (type of research) research''. A complete list of reporting guidelines can be accessed via the equator network: \url{https://www.equator-network.org/}.}
%
%% Only for journal Nursing Reports
%\useofartificialintelligence{Please describe in detail any and all uses of artificial intelligence (AI) or AI-assisted tools used in the preparation of the manuscript. This may include, but is not limited to, language translation, language editing and grammar, or generating text. Alternatively, please state that “AI or AI-assisted tools were not used in drafting any aspect of this manuscript”.}

\acknowledgments{In this section you can acknowledge any support given which is not covered by the author contribution or funding sections. This may include administrative and technical support, or donations in kind (e.g., materials used for experiments). Where GenAI has been used for purposes such as generating text, data, or graphics, or for study design, data collection, analysis, or interpretation of data, please add “During the preparation of this manuscript/study, the author(s) used [tool name, version information] for the purposes of [description of use]. The authors have reviewed and edited the output and take full responsibility for the content of this publication.”}

\conflictsofinterest{Declare conflicts of interest or state ``The authors declare no conflicts of interest.'' Authors must identify and declare any personal circumstances or interest that may be perceived as inappropriately influencing the representation or interpretation of reported research results. Any role of the funders in the design of the study; in the collection, analyses or interpretation of data; in the writing of the manuscript; or in the decision to publish the results must be declared in this section. If there is no role, please state ``The funders had no role in the design of the study; in the collection, analyses, or interpretation of data; in the writing of the manuscript; or in the decision to publish the results''.} 

%%%%%%%%%%%%%%%%%%%%%%%%%%%%%%%%%%%%%%%%%%
%% Optional

%% Only for journal Encyclopedia
%\entrylink{The Link to this entry published on the encyclopedia platform.}

\abbreviations{Abbreviations}{
The following abbreviations are used in this manuscript:
\\

\noindent 
\begin{tabular}{@{}ll}
MDPI & Multidisciplinary Digital Publishing Institute\\
DOAJ & Directory of open access journals\\
TLA & Three letter acronym\\
LD & Linear dichroism
\end{tabular}
}

%%%%%%%%%%%%%%%%%%%%%%%%%%%%%%%%%%%%%%%%%%
%% Optional
\appendixtitles{no} % Leave argument "no" if all appendix headings stay EMPTY (then no dot is printed after "Appendix A"). If the appendix sections contain a heading then change the argument to "yes".
\appendixstart
\appendix
\section[\appendixname~\thesection]{}
\subsection[\appendixname~\thesubsection]{}
The appendix is an optional section that can contain details and data supplemental to the main text---for example, explanations of experimental details that would disrupt the flow of the main text but nonetheless remain crucial to understanding and reproducing the research shown; figures of replicates for experiments of which representative data are shown in the main text can be added here if brief, or as Supplementary Data. Mathematical proofs of results not central to the paper can be added as an appendix.

\begin{table}[H] 
\caption{This is a table caption.\label{tab5}}
%\newcolumntype{C}{>{\centering\arraybackslash}X}
\begin{tabularx}{\textwidth}{CCC}
\toprule
\textbf{Title 1}	& \textbf{Title 2}	& \textbf{Title 3}\\
\midrule
Entry 1		& Data			& Data\\
Entry 2		& Data			& Data\\
\bottomrule
\end{tabularx}
\end{table}

\section[\appendixname~\thesection]{}
All appendix sections must be cited in the main text. In the appendices, Figures, Tables, etc. should be labeled, starting with ``A''---e.g., Figure A1, Figure A2, etc.

%%%%%%%%%%%%%%%%%%%%%%%%%%%%%%%%%%%%%%%%%%
%\isPreprints{}{% This command is only used for ``preprints''.
\begin{adjustwidth}{-\extralength}{0cm}
%} % If the paper is ``preprints'', please uncomment this parenthesis.
%\printendnotes[custom] % Un-comment to print a list of endnotes

\reftitle{References}

% Please provide either the correct journal abbreviation (e.g. according to the “List of Title Word Abbreviations” http://www.issn.org/services/online-services/access-to-the-ltwa/) or the full name of the journal.
% Citations and References in Supplementary files are permitted provided that they also appear in the reference list here. 

%=====================================
% References, variant A: external bibliography
%=====================================
% \bibliography{your_external_BibTeX_file}

%=====================================
% References, variant B: internal bibliography
%=====================================

% ACS format
\isAPAandChicago{}{%
\begin{thebibliography}{999}
% Reference 1
\bibitem[Author1(year)]{ref-journal}
Author~1, T. The title of the cited article. {\em Journal Abbreviation} {\bf 2008}, {\em 10}, 142--149.
% Reference 2
\bibitem[Author2(year)]{ref-book1}
Author~2, L. The title of the cited contribution. In {\em The Book Title}; Editor 1, F., Editor 2, A., Eds.; Publishing House: City, Country, 2007; pp. 32--58.
% Reference 3
\bibitem[Author1 and Author2 (year)]{ref-book2}
Author 1, A.; Author 2, B. \textit{Book Title}, 3rd ed.; Publisher: Publisher Location, Country, 2008; pp. 154--196.
% Reference 4
\bibitem[Author4(year)]{ref-unpublish}
Author 1, A.B.; Author 2, C. Title of Unpublished Work. \textit{Abbreviated Journal Name} year, \textit{phrase indicating stage of publication (submitted; accepted; in press)}.
% Reference 5
\bibitem[Author8(year)]{ref-url}
Title of Site. Available online: URL (accessed on Day Month Year).
% Reference 6
\bibitem[Author6(year)]{ref-proceeding}
Author 1, A.B.; Author 2, C.D.; Author 3, E.F. Title of presentation. In Proceedings of the Name of the Conference, Location of Conference, Country, Date of Conference (Day Month Year); Abstract Number (optional), Pagination (optional).
% Reference 7
\bibitem[Author7(year)]{ref-thesis}
Author 1, A.B. Title of Thesis. Level of Thesis, Degree-Granting University, Location of University, Date of Completion.
\end{thebibliography}
}

% Chicago format (Used for journal: arts, genealogy, histories, humanities, jintelligence, laws, literature, religions, risks, socsci)
\isChicagoStyle{%
\begin{thebibliography}{999}
% Reference 1
\bibitem[Aranceta-Bartrina(1999a)]{ref-journal}
Aranceta-Bartrina, Javier. 1999a. Title of the cited article. \textit{Journal Title} 6: 100--10.
% Reference 2
\bibitem[Aranceta-Bartrina(1999b)]{ref-book1}
Aranceta-Bartrina, Javier. 1999b. Title of the chapter. In \textit{Book Title}, 2nd ed. Edited by Editor 1 and Editor 2. Publication place: Publisher, vol. 3, pp. 54–96.
% Reference 3
\bibitem[Baranwal and Munteanu {[1921]}(1955)]{ref-book2}
Baranwal, Ajay K., and Costea Munteanu. 1955. \textit{Book Title}. Publication place: Publisher, pp. 154--96. First published 1921 (op-tional).
% Reference 4
\bibitem[Berry and Smith(1999)]{ref-thesis}
Berry, Evan, and Amy M. Smith. 1999. Title of Thesis. Level of Thesis, Degree-Granting University, City, Country. Identifi-cation information (if available).
% Reference 5
\bibitem[Cojocaru et al.(1999)]{ref-unpublish}
Cojocaru, Ludmila, Dragos Constatin Sanda, and Eun Kyeong Yun. 1999. Title of Unpublished Work. \textit{Journal Title}, phrase indicating stage of publication.
% Reference 6
\bibitem[Driver et al.(2000)]{ref-proceeding}
Driver, John P., Steffen Rohrs, and Sean Meighoo. 2000. Title of Presentation. In \textit{Title of the Collected Work} (if available). Paper presented at Name of the Conference, Location of Conference, Date of Conference.
% Reference 7
\bibitem[Harwood(2008)]{ref-url}
Harwood, John. 2008. Title of the cited article. Available online: URL (accessed on Day Month Year).
\end{thebibliography}
}{}

% APA format (Used for journal: admsci, behavsci, businesses, econometrics, economies, education, ejihpe, games, humans, ijfs, journalmedia, jrfm, languages, psycholint, publications, tourismhosp, youth)
\isAPAStyle{%
\begin{thebibliography}{999}
\bibitem[Acemoglu et al.(2012)]{ref-acemoglu2012}
Acemoglu, D., Carvalho, V. M., Ozdaglar, A., \& Tahbaz-Salehi, A. (2012). The network origins of aggregate fluctuations. \textit{Econometrica}, \textit{80}(5), 1977--2016.

\bibitem[Alatriste-Contreras(2015)]{ref-alatriste2015}
Alatriste-Contreras, M. G. (2015). The relationship between the key sectors in the European Union economy and the intra-European Union trade. \textit{Journal of Economic Structures}, \textit{4}(14). \url{https://doi.org/10.1186/s40008-015-0024-5}

\bibitem[Alatriste-Contreras \& Lugo(2022)]{ref-alatriste-lugo2022}
Alatriste-Contreras, M. G., \& Lugo, I. (2022). Strategic sectors and the diffusion of the effect of a shock in Mexico for 2008 and 2012.

\bibitem[Alcántara \& Padilla(2020)]{ref-alcantara-padilla2020}
Alcántara, V., \& Padilla, E. (2020). Key sectors in greenhouse gas emissions in Spain: An alternative input--output analysis. \textit{Journal of Industrial Ecology}. \url{https://doi.org/10.1111/jiec.12948}

\bibitem[An et al.(2024)]{ref-an2024}
An, P., Qu, S., Yu, K., \& Xu, M. (2024). Mapping analytical methods between input--output economics and network science. \textit{Journal of Industrial Ecology}, \textit{28}(4). \url{https://doi.org/10.1111/jiec.13493}

\bibitem[André \& Cardenete(2009)]{ref-andre-cardenete2009}
André, F. J., \& Cardenete, M. A. (2009). Designing efficient subsidy policies in a regional economy: A multicriteria decision-making (MCDM)--computable general equilibrium (CGE) approach. \textit{Regional Studies}.

\bibitem[Antoszewski(2019)]{ref-antoszewski2019}
Antoszewski, M. (2019). Assessment of energy-related technological shocks within a CGE model for the Polish economy. \textit{Gospodarka Narodowa}, \textit{1}(297), 9--45. \url{https://doi.org/10.33119/GN/105510}

\bibitem[Armington(1969)]{ref-armington1969}
Armington, P. S. (1969). A theory of demand for products distinguished by place of production. \textit{IMF Staff Papers}, \textit{16}(1), 159--178.

\bibitem[Bahal \& Lenzo(2023)]{ref-bahal-lenzo2023}
Bahal, G., \& Lenzo, D. (2023). Beyond Domar weights: A new measure of systemic importance in production networks. \textit{SSRN Electronic Journal}. \url{https://doi.org/10.2139/ssrn.4512514}

\bibitem[Baqaee \& Farhi(2019a)]{ref-baqaee-farhi2019a}
Baqaee, D. R., \& Farhi, E. (2019a). The macroeconomic impact of microeconomic shocks: Beyond Hulten's theorem. \textit{Econometrica}, \textit{87}(4), 1155--1203. \url{https://doi.org/10.3982/ECTA15202}

\bibitem[Baqaee \& Farhi(2019b)]{ref-baqaee-farhi2019b}
Baqaee, D. R., \& Farhi, E. (2019b). The microeconomic foundations of aggregate production functions. NBER Working Paper No. 25293. \url{https://doi.org/10.3386/w25293}

\bibitem[Baqaee \& Farhi(2020a)]{ref-baqaee-farhi2020a}
Baqaee, D. R., \& Farhi, E. (2020a). Productivity and misallocation in general equilibrium. \textit{Quarterly Journal of Economics}, \textit{135}(1) / NBER WP 24007. \url{https://doi.org/10.3386/w24007}

\bibitem[Baqaee \& Farhi(2020b)]{ref-baqaee-farhi2020b}
Baqaee, D., \& Farhi, E. (2020b). Nonlinear production networks with an application to the Covid-19 crisis. NBER Working Paper No. 27281. \url{https://doi.org/10.3386/w27281}

\bibitem[Bhusal(2025)]{ref-bhusal2025}
Bhusal, L. B. (2025). Key sector identification of Nepal: An integrated approach with linkage and hypothetical extraction method.

\bibitem[Blackburn \& Moreno-Cruz(2021)]{ref-blackburn-moreno2021}
Blackburn, C. J., \& Moreno-Cruz, J. (2021). Energy efficiency in general equilibrium with input-output linkages. \textit{Journal of Environmental Economics and Management} / SSRN. \url{https://doi.org/10.2139/ssrn.3518953}

\bibitem[Bolarinwa(2024)]{ref-bolarinwa2024}
Bolarinwa, T. (2024). Computable general equilibrium (CGE) models: A comprehensive review and future directions. \textit{Management and Economics Review}, \textit{9}(1). \url{https://doi.org/10.24818/mer/2024.01-10}

\bibitem[Burnett et al.(2012)]{ref-burnett2012}
Burnett, P., Cutler, H., \& Davies, S. (2012). Understanding the unique impacts of economic growth variables. \textit{Journal of Regional Science}, \textit{52}(3), 451--468.

\bibitem[Carvalho(2014)]{ref-carvalho2014}
Carvalho, V. M. (2014). From micro to macro via production networks. \textit{Journal of Economic Perspectives}, \textit{28}(4), 23--48.

\bibitem[Carvalho \& Tahbaz-Salehi(2019)]{ref-carvalho-tahbaz2019}
Carvalho, V. M., \& Tahbaz-Salehi, A. (2019). Production networks: A primer. \textit{Annual Review of Economics}, \textit{11}, 635--663.

\bibitem[Chenery \& Watanabe(1958)]{ref-chenery-watanabe1958}
Chenery, H. B., \& Watanabe, T. (1958). International comparisons of the structure of production. \textit{Econometrica}, \textit{26}(4), 487--521.

\bibitem[Çalık et al.(2018)]{ref-calik2018}
Çalık, A., Çizmecioğlu, S., \& Akpınar, A. (2018). An integrated AHP--TOPSIS framework for foreign direct investment in Turkey.

\bibitem[Çırpıcı(2024)]{ref-cirpici2024}
Çırpıcı, Y. A. (2024). Key sector analysis by IO networks: Evidence from Türkiye. \textit{Panoeconomicus}, \textit{71}(3), 395--432. \url{https://doi.org/10.2298/PAN230326023C}

\bibitem[Cutler \& Davies(2010)]{ref-cutler-davies2010}
Cutler, H., \& Davies, S. (2010). The economic consequences of productivity changes: A computable general equilibrium (CGE) analysis. \textit{Regional Studies}, \textit{44}(10), 1415--1426. \url{https://doi.org/10.1080/00343400701654210}

\bibitem[Delzeit et al.(2020)]{ref-delzeit2020}
Delzeit, R., Beach, R., Bibas, R., Britz, W., Chateau, J., Freund, F., \ldots Wojtowicz, K. (2020). Linking global CGE models with sectoral models to generate baseline scenarios: Approaches, opportunities and pitfalls. \textit{Journal of Global Economic Analysis}, \textit{5}(1). \url{https://doi.org/10.21642/jgea.050105af}

\bibitem[DePaolis et al.(2022)]{ref-depaolis2022}
DePaolis, F., Murphy, P., \& De Paolis Kaluza, M. C. (2022). Identifying key sectors in the regional economy: A network analysis approach using input--output data. \textit{Applied Network Science}, \textit{7}, 86. \url{https://doi.org/10.1007/s41109-022-00519-2}

\bibitem[Dervis et al.(1982)]{ref-dervis1982}
Dervis, K., de Melo, J., \& Robinson, S. (1982). \textit{General Equilibrium Models for Development Policy}. Cambridge University Press.

\bibitem[Domar(1961)]{ref-domar1961}
Domar, E. D. (1961). On the measurement of technological change. \textit{Economic Journal}, \textit{71}(284), 709--729.

\bibitem[Ernawati et al.(2024)]{ref-ernawati2024}
Ernawati, Ilyas, \& Asri, M. (2024). Key industrial sectors in the Sulampua area as a result of Nusantara capital city development. \textit{Iqtishaduna}, \textit{13}(1), 50--64.

\bibitem[Fadinger et al.(2022)]{ref-fadinger2022}
Fadinger, H., Ghiglino, C., \& Teteryatnikova, M. (2022). Income differences, productivity, and input-output networks. \textit{American Economic Journal: Macroeconomics}, \textit{14}(2). \url{https://doi.org/10.1257/mac.20180342}

\bibitem[Faridzad(2025)]{ref-faridzad2025}
Faridzad, A. (2025). Applying machine learning to input--output analysis: A new perspective on identifying key Australian industries.

\bibitem[Freeman(1977)]{ref-freeman1977}
Freeman, L. C. (1977). A set of measures of centrality based on betweenness. \textit{Sociometry}, \textit{40}(1), 35--41.

\bibitem[Gabaix(2011)]{ref-gabaix2011}
Gabaix, X. (2011). The granular origins of aggregate fluctuations. \textit{Econometrica}, \textit{79}(3), 733--772.

\bibitem[Guerrieri et al.(2014)]{ref-guerrieri2014}
Guerrieri, L., Henderson, D., \& Kim, J. (2014). Modeling investment-sector efficiency shocks: When does disaggregation matter? \textit{International Economic Review}, \textit{55}(3).

\bibitem[Hirschman(1958)]{ref-hirschman1958}
Hirschman, A. O. (1958). \textit{The Strategy of Economic Development}. Yale University Press.

\bibitem[Hlushko et al.(2025)]{ref-hlushko2025}
Hlushko, A., Laktionov, O., Yanko, A., \& Isaiev, O. (2025). Models for industry differentiation in decision-making systems with an application to the Ukrainian economy. \textit{Modelling and Digitalization}, \textit{37}. \url{https://doi.org/10.32620/reks.2025.3.03}

\bibitem[Hulten(1978)]{ref-hulten1978}
Hulten, C. R. (1978). Growth accounting with intermediate inputs. \textit{Review of Economic Studies}, \textit{45}(3), 511--518.

\bibitem[Ionașcu et al.(2024)]{ref-ionascu2024}
Ionașcu, A. E., Goswami, S. S., Dănilă, A., Horga, M.-G., Barbu, C. A., \& Șerban-Comănescu, A. (2024). Analyzing primary sector selection for economic activity in Romania: An interval-valued fuzzy multi-criteria approach. \textit{Mathematics}, \textit{12}, 1157. \url{https://doi.org/10.3390/math12081157}

\bibitem[İnce(2022)]{ref-ince2022}
İnce, M. R. (2022). Arz, kullanım ve girdi çıktı tabloları arasındaki ilişkiler kapsamında Türkiye ekonomisi için sanayi bazlı girdi çıktı tablosu üretilmesi. \textit{Gazi Journal of Economics and Business}, \textit{8}(1), 145--163.

\bibitem[İnce \& Tarı(2023)]{ref-ince-tari2023}
İnce, M. R., \& Tarı, R. (2023). Verimlilik ve kur şoklarının ihracat, ekonomik büyüme ve refah üzerindeki etkisi: Türkiye üzerine hesaplanabilir genel denge analizi. \textit{Necmettin Erbakan Üniversitesi Siyasal Bilgiler Fakültesi Dergisi}, \textit{5}, 15--32.

\bibitem[İnce \& Tarı(2026)]{ref-ince-tari2026}
İnce, M. R., \& Tarı, R. (2026). Assessing the impact of potential energy wealth on Türkiye's non-energy sectors: A CGE model approach. \textit{Economia Politica}, \textit{43}, 89--127.

\bibitem[Jahangard \& Keshtvarz(2012)]{ref-jahangard-keshtvarz2012}
Jahangard, E., \& Keshtvarz, V. (2012). Identification of key sectors for Iran, South Korea and Turkey economies: A network theory approach. \textit{Iranian Economic Review}, \textit{16}(32).

\bibitem[Jamwal et al.(2021)]{ref-jamwal2021}
Jamwal, A., Agrawal, R., Sharma, M., \& Kumar, V. (2021). Review on multi-criteria decision analysis in sustainable manufacturing decision making. \textit{International Journal of Sustainable Engineering}.

\bibitem[Jiang et al.(2024)]{ref-jiang2024}
Jiang, L., Kim, Y. S., \& Zhang, L. (2024). Sectoral shocks, production linkages, and business cycles in China. \textit{China Economic Review}, \textit{87}, 102211.

\bibitem[Jin \& Li(2025)]{ref-jin-li2025}
Jin, T., \& Li, Z. (2025). Sectoral linkages and their influence on structural change: The case of China. \textit{PLOS ONE}, \textit{20}(9). \url{https://doi.org/10.1371/journal.pone.0330908}

\bibitem[Jin et al.(2022)]{ref-jin2022}
Jin, Y., Xu, Y., Li, R., Zhao, C., \& Yuan, Z. (2022). Comprehensive evaluation of China’s input--output sector status based on the entropy weight--social network analysis method. \textit{Sustainability}, \textit{14}, 14588. \url{https://doi.org/10.3390/su142114588}

\bibitem[Khan(2020)]{ref-khan2020}
Khan, M. A. (2020). Cross sectoral linkages to explain structural transformation in Nepal. \textit{Structural Change and Economic Dynamics}, \textit{52}. \url{https://doi.org/10.1016/j.strueco.2019.11.005}

\bibitem[Kharazi et al.(2025)]{ref-kharazi2025}
Kharazi, A., Hopkinson, P., Zils, M., \& Kobayashi, T. (2025). A dynamic production network approach to model resource productivity shocks from policy interventions in the UK. \textit{International Economics and Economic Policy}, \textit{22}, 63. \url{https://doi.org/10.1007/s10368-025-00690-8}

\bibitem[Kleinberg(1999)]{ref-kleinberg1999}
Kleinberg, J. M. (1999). Authoritative sources in a hyperlinked environment. \textit{Journal of the ACM}, \textit{46}(5), 604--632.

\bibitem[Lee \& Ishiro(2025)]{ref-lee-ishiro2025}
Lee, S., \& Ishiro, T. (2025). Identifying key industries and assessing their economic impacts in Korea's regional economies: A network centrality and input--output approach from 2005 to 2015. \textit{Evolutionary and Institutional Economics Review}. \url{https://doi.org/10.1007/s40844-025-00322-5}

\bibitem[Leontief(1936)]{ref-leontief1936}
Leontief, W. (1936). Quantitative input and output relations in the economic system of the United States. \textit{Review of Economics and Statistics}, \textit{18}(3), 105--125.

\bibitem[Liu(2019)]{ref-liu2019}
Liu, E. (2019). Industrial policies in production networks. \textit{Quarterly Journal of Economics}, \textit{134}(4), 1883--1948. \url{https://doi.org/10.1093/qje/qjz024}

\bibitem[Lofgren et al.(2002)]{ref-lofgren2002}
Lofgren, H., Harris, R. L., \& Robinson, S. (2002). \textit{A Standard Computable General Equilibrium (CGE) Model in GAMS}. IFPRI.

\bibitem[Martinez \& Orrillo(2024)]{ref-martinez-orrillo2024}
Martinez, T. S., \& Orrillo, M. (2024). Industrial policy in production networks: Which sectors to promote in Brazil? \textit{Revista Tempo do Mundo}, \textit{36}. \url{https://doi.org/10.38116/rtm36art5}

\bibitem[Miller \& Blair(2009)]{ref-miller-blair2009}
Miller, R. E., \& Blair, P. D. (2009). \textit{Input-Output Analysis: Foundations and Extensions} (2nd ed.). Cambridge University Press.

\bibitem[Nugroho(2021)]{ref-nugroho2021}
Nugroho, Y. D. (2021). Analysis of input--output table: Identifying leading sectors in Indonesia (case study in 2010, 2016, and 2020).

\bibitem[Odior \& Arinze(2022)]{ref-odior-arinze2022}
Odior, E. S. O., \& Arinze, S. (2022). The concept of computable general equilibrium models. \textit{International Journal of Research in Commerce and Management Studies}, \textit{4}(2). \url{https://doi.org/10.38193/ijrcms.2022.4201}

\bibitem[Ojaleye \& Narayanan(2022)]{ref-ojaleye-narayanan2022}
Ojaleye, D., \& Narayanan, B. (2022). Identification of key sectors in Nigeria: Evidence of backward and forward linkages from input--output analysis. \textit{SocioEconomic Challenges}, \textit{6}(1), 41--62. \url{https://doi.org/10.21272/sec.6(1).41-62.2022}

\bibitem[Page et al.(1999)]{ref-page1999}
Page, L., Brin, S., Motwani, R., \& Winograd, T. (1999). \textit{The PageRank citation ranking: Bringing order to the web}. Stanford InfoLab.

\bibitem[Pereira López et al.(2021)]{ref-pereira-lopez2021}
Pereira López, X., Węgrzyńska, M. A., \& Fernández Fernández, M. (2021). Methodological contribution to the detection of backward linkages between sectors of the economy. \textit{Argumenta Oeconomica}, \textit{1}(46).

\bibitem[Ramírez-Álvarez et al.(2023)]{ref-ramirez-alvarez2023}
Ramírez-Álvarez, J., Chungandro-Carranco, V., Montenegro-Rosero, N., \& Guevara-Rosero, C. (2023). Central industries in the Ecuadorian input--output network. An application of social network analysis. \textit{Networks and Spatial Economics}, \textit{23}. \url{https://doi.org/10.1007/s11067-023-09605-z}

\bibitem[Rasmussen(1956)]{ref-rasmussen1956}
Rasmussen, P. N. (1956). \textit{Studies in Inter-Sectoral Relations}. Einar Harcks Forlag.

\bibitem[Roson \& Sartori(2016)]{ref-roson-sartori2016}
Roson, R., \& Sartori, M. (2016). Input--output linkages and the propagation of domestic productivity shocks: Assessing alternative theories with stochastic simulation. \textit{Economic Systems Research}, \textit{28}(1), 38--54.

\bibitem[Temurshoev \& Oosterhaven(2014)]{ref-temurshoev-oosterhaven2014}
Temurshoev, U., \& Oosterhaven, J. (2014). Analytical and empirical comparison of policy-relevant key sector measures. \textit{Spatial Economic Analysis}, \textit{9}(3), 284--308. \url{https://doi.org/10.1080/17421772.2014.930168}

\bibitem[Tran et al.(2017)]{ref-tran2017}
Tran, T. K., Sato, H., \& Namatame, A. (2017). Key economic sectors and their transitions: Analysis of world input-output network. \textit{Studies in Computational Intelligence} (Robustness in Econometrics), Springer. \url{https://doi.org/10.1007/978-3-319-50742-2\_23}

\bibitem[Tsoulfidis \& Athanasiadis(2022)]{ref-tsoulfidis-athanasiadis2022}
Tsoulfidis, L., \& Athanasiadis, I. (2022). A new method of identifying key industries: A principal component analysis.

\bibitem[Türegün(2022)]{ref-turegun2022}
Türegün, N. (2022). Financial performance evaluation by multi-criteria decision-making techniques.

\bibitem[Wirkierman et al.(2021)]{ref-wirkierman2021}
Wirkierman, A. L., Bianchi, M., \& Torriero, A. (2021). Leontief meets Markov: Sectoral vulnerabilities through circular connectivity. \textit{Networks and Spatial Economics}, \textit{21}. \url{https://doi.org/10.1007/s11067-021-09551-8}
\end{thebibliography}
}{}

% If authors have biography, please use the format below
%\section*{Short Biography of Authors}
%\bio
%{\raisebox{-0.35cm}{\includegraphics[width=3.5cm,height=5.3cm,clip,keepaspectratio]{Definitions/author1.pdf}}}
%{\textbf{Firstname Lastname} Biography of first author}
%
%\bio
%{\raisebox{-0.35cm}{\includegraphics[width=3.5cm,height=5.3cm,clip,keepaspectratio]{Definitions/author2.jpg}}}
%{\textbf{Firstname Lastname} Biography of second author}

% For the MDPI journals use author-date citation, please follow the formatting guidelines on http://www.mdpi.com/authors/references
% To cite two works by the same author: \citeauthor{ref-journal-1a} (\citeyear{ref-journal-1a}, \citeyear{ref-journal-1b}). This produces: Whittaker (1967, 1975)
% To cite two works by the same author with specific pages: \citeauthor{ref-journal-3a} (\citeyear{ref-journal-3a}, p. 328; \citeyear{ref-journal-3b}, p.475). This produces: Wong (1999, p. 328; 2000, p. 475)

%%%%%%%%%%%%%%%%%%%%%%%%%%%%%%%%%%%%%%%%%%
%% for journal Sci
%\reviewreports{\\
%Reviewer 1 comments and authors’ response\\
%Reviewer 2 comments and authors’ response\\
%Reviewer 3 comments and authors’ response
%}
%%%%%%%%%%%%%%%%%%%%%%%%%%%%%%%%%%%%%%%%%%
\PublishersNote{}
%\isPreprints{}{% This command is only used for ``preprints''.
\end{adjustwidth}
%} % If the paper is ``preprints'', please uncomment this parenthesis.
\end{document}

